% Propietario conservado: /Users/ruben/Documents/ChatGPT/jueces y controles/REVISION_ARGUMENTAL_APERTURAS_CIERRES_20260915/II/manuscript_es/sections/18_comparacion_cuantitativa_II.tex
\section{Evaluation of constants and metrological comparison}
\label{ii:sec:ii-contraste-cuantitativo}

The preceding constructions allow three operations to be separated:
producing a quantity, publishing its coordinate in a chart and comparing that
coordinate with a reference. The APP--TRIT--TPK state preserves sheets,
orientation, quotient, carry and memory before these readings. Its
nonadic prolongation publishes the correlated coordinates
\(\pi,\varphi,e,\alpha\); the determinantal reader and regional return
produce the action sections; the marked incidence \(90/120\) and
the conjugate channels produce the Barbero functional; angular
transport subsequently publishes the CKM coordinates. No entry in the external
tables enters into the choice of those operations.

The evaluations preserved in the source documents and their comparators
are brought together here. The digits of the internal outputs express the precision of
numerical publication, not experimental uncertainty. For a value
\(x\) and a nonzero reference \(x_0\), both in the same chart, we use
\[
 \Delta x=x-x_0,\qquad r_x=\frac{x-x_0}{x_0}.
\]
When the reference has asymmetric marginal uncertainties
\(u_-,u_+\), we also define the descriptive residual
\begin{equation}
 z_x=\begin{cases}
 \Delta x/u_+,&\Delta x\geq0,\\
 \Delta x/u_-,&\Delta x<0.
 \end{cases}
 \label{ii:eq:ii-residuo-marginal}
\end{equation}
This normalization does not assume independence between rows and does not
by itself constitute a joint statistical test.

\subsection{Planck constant and reduced constant: two orientations}

In the chart \(\mathcal U_S\mapsto\mathrm{J\,s}\), the sections already
constructed are
\begin{align}
 \hbar_{\rm pre}&=H_5\,10^{-34}\,\mathrm{J\,s},&
 H_5&=1.0545718176461544696\ldots,\\
 \hbar_{\rm ret}&=\eta_{\rm ret}\,10^{-34}\,\mathrm{J\,s},&
 \eta_{\rm ret}&=1.0545718169150390611\ldots,\\
 h_a&=2\pi\hbar_a,&a&\in\{\mathrm{pre},\mathrm{ret}\}.
 \label{ii:eq:ii-cuatro-coordenadas-planck}
\end{align}
The decade comes from \(\Sigma_H=\varphi\alpha^{16}\); the difference between
the two sections comes from regional return. Multiplication by \(2\pi\)
changes action per radian to action per turn, without introducing another
orientation or altering the relative residual.

The SI reference is
\[
 h_{\rm SI}=6.62607015\times10^{-34}\,\mathrm{J\,s},\qquad
 \hbar_{\rm SI}=\frac{h_{\rm SI}}{2\pi}
 =1.0545718176461563912\ldots\times10^{-34}\,\mathrm{J\,s}.
\]
Both expressions are exact. The infinite expansion of \(\hbar_{\rm SI}\)
is not a measurement uncertainty. Fixing the numerical value of
\(h\) belongs to the definition of the SI in effect since 2019
\cite{ii:bipm2018}; comparison takes place after the HMT construction.

\begin{table}[htbp]
\centering\small
\setlength{\tabcolsep}{5pt}
\begin{tabular}{@{}llr@{}}
\toprule
Magnitud & Coordinate in \(10^{-34}\,\mathrm{J\,s}\) & \(r_x\) relative to SI\\
\midrule
\(\hbar_{\rm pre}\) & \(1.0545718176461544696\) & \(-1.822221\times10^{-15}\)\\
\(\hbar_{\rm ret}\) & \(1.0545718169150390611\) & \(-6.932836\times10^{-10}\)\\
\(h_{\rm pre}\) & \(6.62607014999998793\) & \(-1.822221\times10^{-15}\)\\
\(h_{\rm ret}\) & \(6.62607014540625433\) & \(-6.932836\times10^{-10}\)\\
\bottomrule
\end{tabular}
\caption{Action before and after return. The final digits are
evaluations of the published coordinates; the residuals are not experimental
errors of \(h\) or \(\hbar\).}
\label{ii:tab:ii-planck-comparacion}
\end{table}

\begin{proposition}[Invariance of the comparison under a complete turn]
In a common chart,
\[
 \frac{h_a-h_{\rm SI}}{h_{\rm SI}}
 =\frac{\hbar_a-\hbar_{\rm SI}}{\hbar_{\rm SI}}.
\]
\end{proposition}
\begin{proof}
The two terms of the difference in the numerator on the left contain the common factor
\(2\pi\), which cancels against that in the denominator.
\end{proof}
The pre/ret distinction is quantitatively relevant: replacing one with
the other changes the comparison by several orders of magnitude. The regional
corrector must remain visible and must not be confused with numerical rounding.

\subsection{Barbero--Immirzi: internal value and comparison between constructions}

The functional evaluation is reconstructed from its three contributions:
\begin{align*}
 \Delta S_{90}&=2.95169439297457621139847987861\times10^{-4},\\
 \Delta S_{120}&=1.96835112502951720093738706217\times10^{-5},\\
 \pi AC^*/180&=0.270485322934140078465586458790\ldots,\\
 \gamma_{\rm HMT}&=\pi AC^*/180+12\Delta S_{90}-\Delta S_{120}\\
 &=0.27400767269445927474725526077398041940\ldots.
\end{align*}
This composition preserves the oriented weight \(12:-1\), prior to its
decimal evaluation. Its published summands reconstruct the final
coordinate with an absolute residual smaller than \(3\times10^{-31}\).

The Ghosh--Mitra counting is considered as a bibliographic comparison,
not a measurement of the parameter. Their equation (7), written with
\(\lambda=2\pi\gamma_{\rm GM}\), is
\begin{equation}
 \sum_{j\in\{1/2,1,3/2,\ldots\}}(2j+1)
 e^{-2\pi\gamma_{\rm GM}\sqrt{j(j+1)}}=1.
 \label{ii:eq:ii-gm-comparador}
\end{equation}
The reference is A.~Ghosh and P.~Mitra,
\emph{An improved estimate of black hole entropy in the quantum geometry
approach}, \emph{Physics Letters B} \textbf{616} (2005), 114--117,
\href{https://arxiv.org/abs/gr-qc/0411035}{arXiv:gr-qc/0411035},
equations (7)--(8). The root is evaluated here for that specific equation;
puncture counting is not identified with the HMT generator.

\begin{table}[htbp]
\centering\small
\begin{tabular}{@{}lll@{}}
\toprule
Construction & Dimensionless value & Nature\\
\midrule
Functional \(\Gamma_{6\to8}\) & \(0.2740076726944593\) & HMT incidence and return\\
Root of \eqref{ii:eq:ii-gm-comparador} & \(0.2740668582328300\) & Declared bibliographic counting\\
\midrule
\(\gamma_{\rm HMT}-\gamma_{\rm GM}\) & \(-5.918553837\times10^{-5}\) & Difference between constructions\\
\(\gamma_{\rm HMT}/\gamma_{\rm GM}-1\) & \(-2.159529202\times10^{-4}\) & Diferencia relativa\\
\bottomrule
\end{tabular}
\caption{Barbero--Immirzi comparison in the positive orientation. No
experimental uncertainty or statistical significance is assigned to this table.}
\label{ii:tab:ii-barbero-comparacion}
\end{table}

The external root evaluation is reproducible without arbitrarily
truncating its tail. With \(n=2j\), \(N=128\), and
\(r=e^{-0.27\pi}\), for \(\gamma\geq0.27\),
\[
 \sum_{n>N}(n+1)e^{-\pi\gamma\sqrt{n(n+2)}}
 \leq \frac{r^{N+1}[(N+2)-(N+1)r]}{(1-r)^2}<10^{-43}.
\]
Every summand decreases strictly with \(\gamma\). The endpoints
\(0.27\) and \(0.28\) enclose the root; bisection and the tail bound
control the decimals in the table. This is a check of the
bibliographic comparator, not a route for selecting \(\gamma_{\rm HMT}\).

\subsection{CKM angles, amplitudes, and CP violation}

The angular reader already constructed publishes, in degrees,
\begin{equation}
 (\theta_{12},\theta_{23},\theta_{13},\delta)
 =(13.003313589509,\;2.398056157332,\;
 0.213977382244,\;65.676173123554)^\circ.
 \label{ii:eq:ii-ckm-decimales}
\end{equation}
These values arise from \(M_{\rm CKM}(A,C^*/6,\Delta)^T\) and
\(\delta=9A\). Conversion to radians is performed once before
evaluating sines and cosines. The unitary parameterization of the angular
section gives
\[
 |V_{\rm CKM}|\simeq
 \begin{pmatrix}
 0.974350259&0.225005836&0.003734601\\
 0.224873110&0.973489279&0.041841465\\
 0.008582400&0.041121745&0.999117283
 \end{pmatrix},
\]
and the invariant quartet preserves the phase:
\[
 J=c_{12}c_{23}c_{13}^{\,2}s_{12}s_{23}s_{13}\sin\delta
 =3.1189723326632974\times10^{-5}\quad
 \text{in the owner's evaluation.}
\]
This evaluation uses \eqref{ii:eq:ii-jarlskog} with the coordinates of
\eqref{ii:eq:ii-ckm-decimales}. The nine printed moduli are rounded; unitarity belongs
to the complex matrix before rounding, not to a matrix formed
solely from moduli.

The comparator is the three-generation fit in the
\href{https://pdg.lbl.gov/2025/reviews/rpp2025-rev-ckm-matrix.pdf}
{Particle Data Group, 2025 update, \emph{CKM Quark-Mixing Matrix}},
equations (12.27)--(12.28) and the adjacent value of \(J\). To avoid
altering their uncertainties through a nonlinear transformation, the table
retains the observables published by PDG: \(s_{ij}=\sin\theta_{ij}\),
\(\delta\) in radians and \(J\).

\begin{table}[htbp]
\centering\small
\setlength{\tabcolsep}{4pt}
\begin{tabular}{@{}lrrr@{}}
\toprule
Observable & HMT evaluation & PDG 2025 & \(z_x\)\\
\midrule
\(s_{12}\) & \(0.225007404757\) & \(0.22501\pm0.00068\) & \(-0.003817\)\\
\(s_{23}\) & \(0.041841757010\) & \(0.04183^{+0.00079}_{-0.00069}\) & \(0.014882\)\\
\(s_{13}\) & \(0.003734601164\) & \(0.003732^{+0.000090}_{-0.000085}\) & \(0.028902\)\\
\(\delta\) (rad) & \(1.146265461116\) & \(1.147\pm0.026\) & \(-0.028251\)\\
\(10^5J\) & \(3.118972333\) & \(3.12^{+0.13}_{-0.12}\) & \(-0.008564\)\\
\bottomrule
\end{tabular}
\caption{Marginal comparison in the declared parameterization and edition.
The factor \(10^5\) jointly multiplies the value and uncertainty in the
last row; it does not change its normalized residual.}
\label{ii:tab:ii-ckm-comparacion}
\end{table}

The residuals are calculated with the coordinates before rounding
in the table. The PDG fit combines data and unitarity constraints;
its rows are not independent measurements. Therefore,
\(\sum z_x^2\) is not summed and no number of degrees of freedom is assigned to this table.
The result is a reproducible marginal comparison with that edition,
not a measure of model probability. The historical marginal deviations
preserved in the source document are not transferred to this table:
the excerpt does not identify the edition or covariance of that comparison.

The PMNS sector retains the transport \(U=F_\ell^\dagger F_\nu\), the
spectral domains and the phase data developed in the angular
section. A numerical row requires specifying the frames of the register
used and the realization being compared. This table does not assign CKM figures
to the leptonic reader or publish a new PMNS matrix.

\subsection{Boltzmann and universal gravitation: quantity, chart and reference}

The thermal construction already recovered determines the positive transducer
\(k_B:\mathcal L_\Theta\to\mathcal L_E\) and the relation
\begin{equation}
 k_B\Theta_{{\rm clk},a}\log3
 =\frac{h_a}{108t_0}=\frac{\hbar_a}{t_P},
 \qquad t_P=\frac{54}{\pi}t_0.
 \label{ii:eq:ii-kb-comparacion}
\end{equation}
Local information, action and period appear before the
kelvin--joule chart. For positive bases \(u_E,u_\Theta\), if
\(k_B(u_\Theta)=b\,u_E\),
\(\Theta_{{\rm clk},a}=\theta_a u_\Theta\) and
\(E_{P,a}=\epsilon_a u_E\), the equation is equivalent to
\(b\theta_a=\epsilon_a/\log3\). A change
\(u'_E=s_Eu_E\), \(u'_\Theta=s_\Theta u_\Theta\) produces
\[
 b'=\frac{s_\Theta}{s_E}b,\qquad
 \theta'_a=\frac{\theta_a}{s_\Theta},\qquad
 \epsilon'_a=\frac{\epsilon_a}{s_E}.
\]
This preserves the energy product and makes visible the additional
operation that separates the individual coordinates. For the transducer
of subsection~\ref{ii:sec:ii-kb-aut-desarrollo}, the kelvin--joule chart
used here expresses the SI coordinate
\[
 k_B^{\rm SI}=1.380649\times10^{-23}\,\mathrm{J\,K^{-1}},
\]
the exact value of the SI definition, not an estimate with measurement error.
The function of the transducer and its relation to action have content
of their own; citing its SI coordinate does not constitute an independent numerical
check of an internal choice of both bases.

In gravity, the circular realization fixes
\begin{equation}
 G_a=\vartheta_{108}^{\circ}\frac{c_{\rm int}^{5}t_0^2}{\hbar_a},
 \qquad
 \vartheta_{108}^{\circ}=(54/\pi)^2
 =295.4525715010569\ldots.
 \label{ii:eq:ii-g-comparacion}
\end{equation}
The selector and reciprocity \(G_a\hbar_a=
\vartheta_{108}^{\circ}c_{\rm int}^{5}t_0^2\) are internal quantities of
the declared realization. A dimensional evaluation of \(G_a\) must
preserve, in addition to the orientation \(a\), the time,
velocity and action charts. In the dimensional chart used here, the evaluated value is
\(6.67430\times10^{-11}\,\mathrm{m^3\,kg^{-1}\,s^{-2}}\).
The table distinguishes that coordinate from the
recommended value with its uncertainty.

\begin{table}[htbp]
\centering\small
\begin{tabularx}{\linewidth}{@{}lXX@{}}
\toprule
Objeto & Internal output or relation & Referencia dimensional\\
\midrule
\(k_B\) & Positive transducer; \(k_B\Theta_{{\rm clk},a}=E_{P,a}/\log3\).
& \(1.380649\times10^{-23}\,\mathrm{J/K}\), exact in SI.\\[4pt]
\(G_a\) & Selector \((54/\pi)^2\); \(G_a=\vartheta_{108}^{\circ}c_{\rm int}^5t_0^2/\hbar_a\).
& \(6.67430(15)\times10^{-11}\,\mathrm{m^3/(kg\,s^2)}\), CODATA 2022.\\
\bottomrule
\end{tabularx}
\caption{Internal normalization and chart reference. The standard relative
uncertainty of the reference for \(G\) is approximately
\(2.24743\times10^{-5}\); \(k_B\) has no definitional uncertainty.
No predictive residual is calculated from a coordinate that has already
been identified as a chart reference.}
\label{ii:tab:ii-kb-g-cartas}
\end{table}

The external sources for this last table are the
\href{https://www.bipm.org/en/measurement-units/si-defining-constants}
{BIPM defining constants} and the
\href{https://physics.nist.gov/cuu/Constants/Table/allascii.txt}
{NIST CODATA 2022 table} \cite{ii:codata2022}.
The experimental uncertainty of \(G\) is not transferred to \(h\),
\(\hbar\) or \(k_B\). Nor does equality of digits in a chart
remove the need to identify the reader that produces the quantity.

\subsection{Reproduction of the evaluations}

The program \texttt{verificar\_comparaciones\_II.py}, included with the
sources, receives the internal coordinates published in files separate
from the external data. It checks the composition of the Barbero functional,
the conversion between \(h\) and \(\hbar\), the bibliographic root and its tail,
the CKM moduli, the quartet \(J\) and each printed residual. It uses
75-digit decimal arithmetic and elementary functions evaluated after
freezing the internal outputs. Execution instructions
and program versions are included in the reproducible package.
The working precision does not increase the digits available from the
source document: only digits stable under its rounding are published.
The check records the files used and their hashes, so that
a change in evaluation or external edition remains visible.

% NUMCHECK hbar_SI_coefficient=1.0545718176461563912 tol=1e-19
% NUMCHECK h_pre_coefficient=6.62607014999998793 tol=6e-18
% NUMCHECK h_ret_coefficient=6.62607014540625433 tol=6e-18
% NUMCHECK relative_pre=-1.822221e-15 tol=6e-22
% NUMCHECK relative_ret=-6.932836e-10 tol=6e-17
% NUMCHECK gamma_GM=0.2740668582328300 tol=6e-17
% NUMCHECK gamma_difference=-5.918553837e-5 tol=6e-15
% NUMCHECK gamma_relative=-2.159529202e-4 tol=6e-14
% NUMCHECK ckm.s12.value=0.225007404757 tol=6e-13
% NUMCHECK ckm.s23.value=0.041841757010 tol=6e-13
% NUMCHECK ckm.s13.value=0.003734601164 tol=6e-13
% NUMCHECK ckm.delta_rad.value=1.146265461116 tol=6e-13
% NUMCHECK ckm.s12.marginal_residual=-0.003817 tol=6e-7
% NUMCHECK ckm.s23.marginal_residual=0.014882 tol=6e-7
% NUMCHECK ckm.s13.marginal_residual=0.028902 tol=6e-7
% NUMCHECK ckm.delta_rad.marginal_residual=-0.028251 tol=6e-7
% NUMCHECK ckm.J.marginal_residual=-0.008564 tol=6e-7
% NUMCHECK gravity_selector=295.4525715010569 tol=6e-14
% NUMCHECK G_relative_reference_uncertainty=2.24743e-5 tol=6e-11
